\documentclass[a4paper,12pt]{article}

% ----------------------------
% Pacchetti utili
% ----------------------------
\usepackage[utf8]{inputenc}
\usepackage[T1]{fontenc}
\usepackage[italian]{babel}
\usepackage{graphicx}
\usepackage{tabularx}
\usepackage[table]{xcolor}
\definecolor{lightblue}{RGB}{225,240,255}
\definecolor{gold}{RGB}{255,215,0}
\usepackage{geometry}
\usepackage{setspace}
\usepackage{calc}
\usepackage{array}
\usepackage{fancyhdr}
\usepackage{comment}
\usepackage{tikz}
\usepackage{float}
\usepackage{amsmath}
\usepackage{pgf-pie}
\usepackage{longtable}
\usepackage[colorlinks=true, linkcolor=blue, urlcolor=blue, citecolor=blue]{hyperref}
\usepackage{pgffor}

% ----------------------------
% Impostazioni pagina
% ----------------------------
\geometry{
    top=2cm,
    bottom=2cm,
    left=2cm,
    right=2cm
}

\setstretch{1.2}

% ----------------------------
% Dati personalizzabili
% ----------------------------
\newcommand{\Gruppo}{Atlas}
\newcommand{\Email}{\href{mailto:team9.atlas@gmail.com}{\textcolor{blue}{\underline{team9.atlas@gmail.com}}}}
\newcommand{\TitoloDocumento}{Manuale Utente}
\newcommand{\DataUltimaModifica}{2026/03/15}
\newcommand{\LogoGruppo}{../../../Assets/AtlasLogo.png}

% --- Nuove variabili aggiunte ---
\newcommand{\VersioneDocumento}{0.1.0}
\newcommand{\TipoDocumento}{Esterno} 

\pagestyle{fancy}
\fancyhf{}
\fancyhead[L]{\Gruppo}
\fancyhead[R]{Documento: \TipoDocumento}
\fancyfoot[C]{\thepage}

% larghezza della colonna dei nomi
\newlength{\namecol}
\setlength{\namecol}{4.5cm}
\setlength{\headheight}{15pt} % messo per evitare warning di fancyhdr


\newlength{\colw}
\setlength{\colw}{1.5cm}

\setcounter{secnumdepth}{5}
\setcounter{tocdepth}{5}

\definecolor{mioverde}{RGB}{20,150,60}

% ----------------------------
% Inizio documento
% ----------------------------
\begin{document}

% ----------------------------
% Prima pagina
% ----------------------------
\begin{titlepage}

    \begin{center}

        % Logo
        \vspace*{0cm}
        \begin{tikzpicture}
            \clip (0,-0.1) circle (5.6cm);
            \node at (0,0) {\includegraphics[width=12cm]{\LogoGruppo}};
        \end{tikzpicture}\\[0.8cm]

        % Barra superiore
        \noindent\rule{\textwidth}{0.4pt}

        % Titolo
        \vspace{1cm}
        {\Huge \textbf{\TitoloDocumento}}\\[0.4cm]
        {\large Progetto di Ingegneria del Software A.A. 2025/2026}\\[0.8cm]
        {\large Versione: \VersioneDocumento}
        \vspace{1cm}

        % Barra inferiore
        \noindent\rule{\textwidth}{0.4pt}

    \end{center}

    % Informazioni in basso
    \vfill
    \noindent
    \begin{minipage}{0.5\textwidth}
        \raggedright
        \textbf{Autore:} \Gruppo\\
        \textbf{Ultima modifica:} \DataUltimaModifica
    \end{minipage}%
    \begin{minipage}{0.5\textwidth}
        \raggedleft
        \textbf{Tipo di documento:} \TipoDocumento
    \end{minipage}

\end{titlepage}



\section*{Registro delle modifiche}
\begin{center}
    \rowcolors{2}{lightblue}{white}
    \begin{tabularx}{\textwidth}{|l|l|l|X|X|}
        \hline
        \rowcolor{lightgray}
        \textbf{Versione}  & \textbf{Data}       & \textbf{Autore}      & \textbf{Verificatore} & \textbf{Descrizione}  \\
        \hline
        0.1.0 & 2025/03/15 & Bilal Sabic & Federico Simonetto & Prima stesura documento \\
        \hline
    \end{tabularx}
\end{center}

\newpage

\tableofcontents

\newpage

\listoffigures

%\newpage

%\listoftables

\newpage

\section{Introduzione}

    \subsection{Scopo del documento}
        Il presente documento ha l'obiettivo di fornire una guida all'utilizzo dell'applicazione sviluppata dal team \textit{Atlas}, realizzata su richiesta del proponente \textit{Bluewind S.r.l.} nell'ambito del capitolato \textbf{C1 - Automated EN18031 Compliance Verification}. \newline
        Il documento è finalizzato a supportare l'utente finale nelle fasi di installazione e di utilizzo dell'applicazione software. A tal fine vengono descritti i requisiti hardware e software necessari, nonché le procedure da seguire per l'installazione e per il corretto utilizzo del sistema.

    \subsection{Scopo del prodotto}
        L'obiettivo del prodotto è automatizzare il processo di verifica della conformità degli apparecchi che utilizzano onde radio rispetto alla normativa europea EN18031, definita dalla Direttiva RED. \newline
        L'applicazione sviluppata consente di effettuare tale verifica tramite l'esecuzione di specifici \textit{decision tree} associati ai diversi asset del dispositivo. Al termine dell'esecuzione, il sistema fornisce un risultato chiaro della valutazione, espresso attraverso uno dei seguenti esiti: \textit{PASS}, \textit{FAIL} oppure \textit{NON APPLICABLE}.

    \subsection{Contesto d'uso}
        Come già esplicitato nel documento \href{https://atlasteam9.github.io/Atlas/docs/PB/documenti/esterni/Analisi%20dei%20Requisiti.pdf}{Analisi dei Requisiti v2.0.0}, gli utenti finali dell'applicazione \textbf{Automated EN18031 Compliance Verification} appartengono alla categoria degli "utenti esperti". Essi possiedono pertanto competenze tecniche in ambito di sicurezza informatica e conoscenze della normativa europea a cui il capitolato fa riferimento. \newline
        Per tale motivo, nel presente documento viene fatto deliberatamente uso di terminologia tecnica.

    % T. nell'ultimo DdB ha detto che il Manuale, essendo diretto ad utenti esterni, (contrariamente al glossario che è interno) ha bisogno del suo specifico glossario
    
    \subsection{Glossario}
        All'interno della documentazione prodotta dal team possono comparire termini suscettibili di incomprensioni o ambiguità. Per evitare questo, è disponibile un glossario
        contenente i termini tecnici e le loro definizioni. Un termine è consultabile nel glossario se è indicato con la notazione \textit{parola\textsubscript{\href{https://atlasteam9.github.io/Atlas/glossario.html}{G}}}.
        Premendo sulla G a pedice, l'utente verrà indirizzato alla pagina web del \textit{Glossario v.2.0.0}.
    

    \subsection{Riferimenti}
        \subsubsection{Riferimenti normativi}
            \begin{itemize}
                \item Riferimento al capitolato 1 dell'azienda proponente:\newline \textbf{Bluewind S.r.l - Automated EN18031 Compliance Verification}\newline \url{https://www.math.unipd.it/~tullio/IS-1/2025/Progetto/C1.pdf}
                \item Riferimento al Way of Working del team Atlas: \newline \textbf{Norme di Progetto v.2.0.0}\newline \url{da_mettere_una_volta_scritto}
            \end{itemize}

\newpage
\section{Requisiti}
    \subsection{Requisiti hardware}

    \subsection{Requisiti software}
        L'applicazione richiede l'installazione dei seguenti software:
        \begin{itemize}
            \item \textbf{Docker Engine};
            \item \textbf{Git}.
        \end{itemize}
        Il prodotto software viene distribuito tramite container Docker; pertanto non è necessario installare librerie o servizi di supporto aggiuntivi. \newline
        \textbf{Git} è necessario per clonare la repository GitHub contenente il progetto, mentre \textbf{Docker} consente di eseguire l'applicazione all'interno di un container, garantendo un ambiente di esecuzione isolato e configurato correttamente.
\newpage

\section{Installazione}

\newpage

\section{Istruzioni per l'uso}


\end{document}            